% =========================================================
\chapter{Lecture 5 : Applied cryptography}
% =========================================================

\begin{rembox}[title=In a nutshell]
\textbf{What is this lecture about?} How Alice can send a message to Bob so that Eve,
who sees everything on the wire, \textbf{cannot read it}. We start from broken historical
ciphers (Caesar, substitution) to learn how attackers think, define the different kinds
of attackers, see the only \emph{perfectly} secure cipher (one-time pad) and why it is
impractical, then how real systems approximate it (\textbf{stream ciphers}), and finally
how two strangers can agree on a secret key over a public channel
(\textbf{Diffie-Hellman}). \\
\textbf{Why is it useful?} Cryptography protects every HTTPS page, every VPN, every
encrypted backup and every messaging app. Understanding its basics tells you
\emph{what} a cipher guarantees (and what it does not, e.g.\ integrity), and why the
golden rule is: \textbf{never design your own crypto, use vetted standards}.
\end{rembox}

\section{Module 1 -- Introduction}

\begin{rembox}[title=Why cryptography matters]
\begin{itemize}
  \item \textbf{Data in transit} --- securing communication (e.g.\ HTTPS, VPNs).
  \item \textbf{Data at rest} --- securing stored information (e.g.\ encrypted
        backups).
\end{itemize}
Cryptography is also the building block for \textbf{integrity} (ensuring data has not
been tampered with: hashing, signatures), \textbf{authentication} (verifying identity:
protocols), \textbf{anonymity}, etc. This lecture focuses on the basics of using
cryptography for \textbf{confidentiality}.
\end{rembox}

\begin{defnbox}[title=Secure communication over an insecure channel]
The core problem: we need a way for \textbf{Alice} to send a message to \textbf{Bob},
so that \textbf{Eve} (the adversary) \textbf{cannot read it}. \\
Confidentiality for data at rest can be seen as a particular case, where
Alice $=$ Bob (you ``send'' the data to your future self).
\end{defnbox}

\begin{defnbox}[title=Encryption and decryption]
\textbf{Confidentiality}: information cannot be accessed by unauthorized parties.
\begin{itemize}
  \item \textbf{Encryption algorithm} --- cryptographic instructions that convert a
        \textbf{plaintext} message into \textbf{ciphertext}.
  \item \textbf{Decryption algorithm} --- cryptographic instructions that convert a
        ciphertext message into plaintext.
\end{itemize}
Both encryption and decryption algorithms require a \textbf{key}: encryption cannot be
reversed without the key. As functions, with plaintext $M$ (original message), key $K$
(secret value) and ciphertext $C$ (scrambled message):
\[
  C = E_K(M) \qquad\text{and}\qquad M = D_K(C).
\]
\end{defnbox}

\begin{algbox}[title=Cryptographic algorithms for confidentiality]
\begin{enumerate}
  \item Get a key $k$ shared between Alice and Bob.
  \item Alice encrypts the message: $m \to \mathrm{Enc}(k, m)$.
  \item Alice sends $\mathrm{Enc}(k, m)$ over the insecure channel.
  \item Bob decrypts: $\mathrm{Dec}(k, \mathrm{Enc}(k, m)) \to m$.
\end{enumerate}
\end{algbox}

\begin{rembox}[title=Intuition]
Think of the algorithm as a \textbf{lock design} and the key as \textbf{the key}: the
lock design is public (open design, Lecture 2 / Kerckhoffs), only the key is secret.
Step 1 (sharing the key) is the hard part, and comes back at the end of the lecture.
\end{rembox}

\begin{defnbox}[title=Invertibility (the core requirement)]
The decryption function $D_K$ \textbf{must} be the \textbf{inverse} of $E_K$:
\[
  \forall K, M,\quad D_K\big(E_K(M)\big) = M .
\]
If not, Bob cannot recover the original message $M$. \\
\textbf{Security requirement}: to provide security, these functions must also be
\textbf{hard to invert without knowing the key} $K$.
\end{defnbox}

\begin{rembox}[title=Two requirements that pull in opposite directions]
\textbf{Correctness}: \emph{with} the key, inverting is easy. \textbf{Security}:
\emph{without} the key, inverting is hard. The whole art of cryptography is to get both
at the same time.
\end{rembox}

\begin{defnbox}[title=Hard to invert: the key space and brute force]
\textbf{Key space}: the set of all possible keys. \\
\textbf{Concrete meaning of ``hard to invert''}: to find the plaintext $M$, in an ideal
cryptosystem, the only viable attack for Eve must be to \textbf{try every single possible
secret key} $K$ (a \textbf{brute-force attack}). \\
\textbf{Goal}: a secure scheme forces the attacker into this situation.
\end{defnbox}

\begin{defnbox}[title=Caesar's cipher]
\begin{itemize}
  \item \textbf{Encryption}: shift each letter by a fixed number $K$.
  \item \textbf{Decryption}: shift each letter by $-K$.
  \item \textbf{Key space}: only 25 possible keys.
\end{itemize}
\end{defnbox}

\begin{exbox}[title={Caesar with $K = 3$}]
\texttt{ATTACK AT DAWN} $\xrightarrow{K=3}$ \texttt{DWWDFN DW GDZQ}
($A \to D$, $T \to W$, \ldots). \\
\textbf{Security insight}: $\log_2(25) \approx 4.6$ bits of security --- \textbf{far too
small} for real-world use (a brute force is trivial). It is still pedagogically useful to
introduce a first real attack.
\end{exbox}

\begin{defnbox}[title=Frequency analysis attack]
\begin{itemize}
  \item Eve knows that Alice wrote a long-enough message for Bob in English.
  \item She knows that the most frequent letter in English is `e'.
  \item She identifies the most frequent letter in the ciphertext: it likely corresponds
        to `e'.
  \item She then directly recovers the key (a number between 0 and 25).
\end{itemize}
\end{defnbox}

\begin{rembox}[title=Why it works]
Caesar maps every `e' to the \emph{same} ciphertext letter, so the ciphertext keeps the
\textbf{statistics} of English. Eve does not even need to try the 25 keys: one look at the
letter counts gives the shift.
\end{rembox}

\begin{defnbox}[title=Substitution cipher]
Encryption/decryption by \textbf{permuting the alphabet}: each letter is mapped to a
unique, different letter, defined by a \textbf{permutation} of the 26 letters (the key).
\begin{center}
\texttt{ABCDEFGHIJKLMNOPQRSTUVWXYZ} \\
\texttt{HOWBUGIACRYEVZXPJQMSNTFDKL}
\end{center}
\textbf{Key space}: $26! \approx 4.03 \times 10^{26}$ possible keys, i.e.\
$\log_2(26!) \approx 88.4$ bits of security?
\end{defnbox}

\begin{warnbox}[title={The flaw, again}]
This key space looks big enough, but the cipher is \textbf{still easily broken} by the
same frequency analysis. \textbf{Vulnerability}: the cipher \textbf{preserves the
statistical properties} of the plaintext language (e.g.\ single-letter frequencies). With
a long enough ciphertext, statistical analysis breaks the cipher, rendering the full $26!$
key space \textbf{irrelevant}. \\
E.g.\ on the whole Harry Potter series encrypted with one key, Eve guesses the most
likely letter mappings, uses a dictionary to rule out keys that give non-English text,
and finds $K$ in a few tries.
\end{warnbox}

\begin{rembox}[title=Goal of the adversary]
Eve does not necessarily need the \emph{full} key: depending on the scenario, recovering
part of it (the mapping of some letters) may be enough. Sometimes the goal is even
smaller: Eve may only want to know whether Alice sent ``Yes'' or ``No''.
\end{rembox}

\begin{keybox}[title=Lesson]
A \textbf{big key space is necessary but not sufficient}. Security is about the
\emph{best} attack, not only brute force.
\end{keybox}

\begin{defnbox}[title=Bits of security]
\textbf{Goal}: an $N$-bit key should offer security as close to $N$ bits as possible,
i.e.\ require $2^N$ attempts. \\
\textbf{The break}: if a 1024-bit key is broken with only $2^{40}$ attempts (40 bits of
security), the algorithm is considered \textbf{broken}. \\
\textbf{Caesar lesson}: the 25-key space was cleverly reduced to \textbf{1 possibility};
similarly the $\approx 88$ bits of the substitution cipher were reduced to a few tries.
\end{defnbox}

\begin{rembox}[title=Orders of magnitude]
Each extra bit \textbf{doubles} the work. $2^{40} \approx 10^{12}$ operations is easy for
a laptop; $2^{128}$ is far beyond all computers on Earth for billions of years. ``Broken''
does not mean ``practically decrypted today'', it means \textbf{faster than brute
force}.
\end{rembox}

\section{Module 2 -- Adversaries in cryptography}

\begin{rembox}[title=Back to threat models]
As in Lecture 1, security is always measured \textbf{relative to the power of the
adversary}. Against Caesar and substitution, Eve's capabilities were very limited: she
observed a long-enough ciphertext of English text (and had some compute power to test a
few keys). This is not the only possible adversary model.
\end{rembox}

\begin{defnbox}[title=Types of adversaries]
\begin{itemize}
  \item \textbf{Passive eavesdropper} --- Eve can only \emph{read} the ciphertext.
  \item \textbf{Active attacker} --- Eve can somehow \textbf{influence the system},
        e.g.\ convince Alice to send Bob specific messages, observe the precise time
        Bob takes to run his algorithm, etc.: anything that goes beyond just observing
        communication.
\end{itemize}
\end{defnbox}

\begin{defnbox}[title=Known Plaintext Attack (KPA)]
\textbf{Scenario}: Eve is given access to multiple pairs $(M, C)$ of messages and their
corresponding ciphertexts, all encrypted with the secret key $K$. \\
\textbf{Goal}: from these pairs, guess the key $K$ to decrypt other messages.
\end{defnbox}

\begin{propbox}[title=Why is KPA a realistic model?]
\begin{itemize}
  \item \textbf{Standard headers} --- many encrypted messages contain predictable parts
        (``From:'', timestamps, packet destinations, \ldots).
  \item \textbf{Malicious guessing} --- if Eve guesses a message (or part of it) and
        observes its encryption, she has an $(M, C)$ pair.
\end{itemize}
\end{propbox}

\begin{exbox}[title=KPA and the substitution cipher]
Under KPA the substitution cipher becomes \textbf{even more broken}: if Eve knows that
$M = $ \texttt{THE QUICK BROWN FOX}, she immediately recovers the mapping of many
letters, drastically reducing the key space without any statistical analysis.
\end{exbox}

\begin{defnbox}[title=Chosen Plaintext Attack (CPA)]
\textbf{Scenario}: Eve can convince Alice (or the system) to encrypt \textbf{arbitrary
messages chosen by Eve} with the secret key. \\
\textbf{More abstractly}: for any chosen plaintext $m$, Eve gets $E_K(m)$. We say that
Eve has access to an \textbf{encryption oracle}.
\end{defnbox}

\begin{rembox}[title=KPA vs.\ CPA]
In KPA, Eve does not choose which pairs she sees; in CPA she \emph{chooses} $m$ and gets
$(m, \mathrm{Enc}(k, m))$. Hence CPA is a \textbf{strictly stronger} model than KPA:
\[
  \text{passive (ciphertext only)} \;<\; \text{KPA} \;<\; \text{CPA}.
\]
A scheme secure against a stronger adversary is also secure against the weaker ones.
\end{rembox}

\begin{propbox}[title=Why is CPA a realistic model?]
\begin{itemize}
  \item \textbf{Messaging app} --- with a few minutes of access to Alice's phone
        (Signal, WhatsApp), Eve can send \emph{any} data she chooses to Bob and obtain an
        $(M, C)$ pair.
  \item \textbf{Encryption API} --- Eve interacts with an endpoint where she provides
        data and receives the encrypted output.
\end{itemize}
\end{propbox}

\begin{exbox}[title=CPA and the substitution cipher]
The substitution cipher is \textbf{completely broken} by a CPA: Eve asks for the
encryption of $M = $ \texttt{abcdefghijklmnopqrstuvwxyz}; the ciphertext \emph{is} the
entire key (the full permutation), in just \textbf{one} query!
\end{exbox}

\begin{defnbox}[title=Side-channel attacks]
Eve not only accesses an ideal oracle, but also \textbf{measures physical
characteristics of the actual implementation} while it encrypts/decrypts.
\begin{itemize}
  \item \textbf{Timing attack} --- Eve measures \textbf{how long} it takes to encrypt a
        given message or decrypt a ciphertext.
  \item \textbf{Power analysis} --- Eve observes the \textbf{energy consumed} by the
        device (e.g.\ a credit card) doing the crypto.
\end{itemize}
\end{defnbox}

\begin{propbox}[title=Why are side channels realistic?]
\begin{itemize}
  \item \textbf{Malicious ownership} --- you often own a device (phone, smart card,
        embedded chip) holding keys of \emph{another} principal (Google, your bank) that
        are meant to be unreadable to you.
  \item \textbf{DRM} --- the decryption key of a movie/song is embedded in your media
        player or console; by measuring power or timing, an attacker could extract it.
\end{itemize}
\end{propbox}

\begin{rembox}[title=Algorithm vs.\ implementation]
A mathematically perfect cipher can still leak through its \textbf{implementation}
(e.g.\ code that takes longer when a key bit is 1). Side channels attack the code and
the hardware, not the math --- a reminder that security is about the whole system. As
always, new adversary models can be defined, but always as a \textbf{strategic adversary
with limited, well-defined resources}.
\end{rembox}

\section{Module 3 -- Perfect secrecy and the one-time pad}

\begin{rembox}[title={Reminder: XOR ($\oplus$)}]
\begin{center}
\begin{tabular}{cc|c}
$a$ & $b$ & $a \oplus b$ \\ \hline
0 & 0 & 0 \\ 0 & 1 & 1 \\ 1 & 0 & 1 \\ 1 & 1 & 0
\end{tabular}
\end{center}
Key properties: $a \oplus a = 0$, $a \oplus 0 = a$, hence $a \oplus b \oplus b = a$
(XOR-ing twice with the same value cancels out). XOR with a random bit flips the input
bit with probability $1/2$: the output is itself random.
\end{rembox}

\begin{defnbox}[title=One-Time Pad (OTP)]
How to remove frequency analysis? Make sure that letters are \textbf{not} encrypted to
the same values. \\
\textbf{Key}: a string $k$ of \textbf{random} bits, \textbf{as long as the message},
pre-shared between Alice and Bob.
\[
  \mathrm{Enc}(k, m) = m \oplus k \qquad\qquad \mathrm{Dec}(k, c) = c \oplus k
\]
Decryption works because $(m \oplus k) \oplus k = m$.
\end{defnbox}

\begin{exbox}[title=OTP on ``YEAH'']
\begin{tabular}{ll}
Message (ASCII \texttt{59454148}) & \texttt{01011001010001010100000101001000} \\
$\oplus$ OTP key                   & \texttt{01110101000111010100101001001010} \\ \hline
Ciphertext                         & \texttt{00101100010110000000101100000010}
\end{tabular} \\[4pt]
XOR-ing the ciphertext with the same key gives back the message. Encrypting ``YEAH''
twice with \emph{different} keys gives \emph{different} ciphertexts, and ``YEAH'' and
``NOPE'' can produce any ciphertext depending on the key: the ciphertext alone tells
nothing.
\end{exbox}

\begin{defnbox}[title=Perfect secrecy]
For any ciphertext $C$, \textbf{every possible plaintext $M$ is equally likely} as
before seeing $C$: the ciphertext gives the adversary \textbf{zero new information} about
$M$. Formally:
\[
  \forall m, c,\quad \PR\big(M = m \mid E_K(M) = c\big) = \PR(M = m).
\]
\end{defnbox}

\begin{proofbox}[title=The OTP is perfectly secure]
Messages $m$ have some prior probability $\PR(M = m)$ (``Yes'' may be more likely than
``Autobus''), and the key $K$ is uniform on $n$-bit strings, independent of $M$. Since
$E_K(m) = c \iff K = m \oplus c$:
\begin{align*}
  \PR(M = m \mid C = c)
  &= \frac{\PR(M = m \cap K = m \oplus c)}{\PR(C = c)}
   = \frac{\PR(M = m)\,\PR(K = m \oplus c)}{\PR(C = c)} \\
  &= \frac{\PR(M = m)\, 2^{-n}}{2^{-n}}
   = \PR(M = m),
\end{align*}
using independence of $M$ and $K$, and
$\PR(C = c) = \sum_{m'} \PR(M = m')\,\PR(K = m' \oplus c) = 2^{-n}$.
\end{proofbox}

\begin{rembox}[title=Intuition]
For a given ciphertext $c$, \emph{every} message $m$ is possible: there is exactly one key
($k = m \oplus c$) that maps it to $c$, and all keys are equally likely. So even with
infinite computing power, Eve cannot tell which message was sent.
\end{rembox}

\begin{defnbox}[title={One-time: never reuse the key}]
One-time pads are called like this because they must \textbf{not} be reused. If the same
key encrypts two messages:
\[
  C_1 \oplus C_2 = (M_1 \oplus K) \oplus (M_2 \oplus K) = M_1 \oplus M_2 ,
\]
the key cancels out. Eve cannot recover the full messages directly, but she learns
information about the plaintexts, which can feed frequency analysis and help recover
them.
\end{defnbox}

\begin{exbox}[title=Spaces vs.\ letters in ASCII]
In ASCII, spaces start with \texttt{00} and letters with \texttt{01}. In $M_1 \oplus M_2$:
\begin{itemize}
  \item 2 letters: $01 \oplus 01 = 00$;
  \item 2 spaces: $00 \oplus 00 = 00$;
  \item letter $+$ space: $01 \oplus 00 = 01$ $\Rightarrow$ one of the messages has a
        space at this position!
\end{itemize}
\end{exbox}

\begin{exbox}[title={Exercise --- Key reuse}]
Suppose Alice reuses the key $K$ to encrypt \textbf{two single-word English messages of
length 5}, $M_1$ and $M_2$, producing $C_1 = M_1 \oplus K$ and $C_2 = M_2 \oplus K$. \\
\textbf{Task}: show that Eve can ``likely'' recover the key and then the two messages
$M_1$ and $M_2$. (Hint: Eve can look up an English dictionary of all the words of
length 5.)
\end{exbox}

\begin{defnbox}[title=The key problem of the OTP]
\textbf{Practical limitation}: OTP requires keys that are \textbf{truly random,
message-length, and used only once}. \\
\textbf{The showstopper}: securely pre-distributing message-length keys for real-world
applications is \textbf{impractical}.
\end{defnbox}

\begin{rembox}[title=Why is it a showstopper?]
To send 1~GB securely you first need to share 1~GB of secret randomness \emph{securely}
--- but if you had a secure way to share 1~GB, you could just send the message that way!
And reusing keys is not an option (see above).
\end{rembox}

\begin{defnbox}[title={OTP: integrity is NOT included (malleability)}]
The OTP guarantees \textbf{confidentiality} (nobody can read it), but it \textbf{does
not} guarantee \textbf{integrity} (nobody can tamper with it).
\end{defnbox}

\begin{exbox}[title=Bit-flipping attack]
Eve flips bits of $C$ with a mask $\Delta$ (e.g.\ only the first bit set):
$C' = C \oplus \Delta$. Bob decrypts
\[
  M' = C' \oplus K = (M \oplus K \oplus \Delta) \oplus K = M \oplus \Delta .
\]
Eve successfully changed the message ($M'$ is $M$ with one bit flipped)
\textbf{without knowing $K$}. If Eve knows the format (e.g.\ ``PAY 0100 CHF''), she can
change the amount in a controlled way.
\end{exbox}

\begin{warnbox}[title={Confidentiality $\neq$ integrity}]
Encryption hides the content; it does \textbf{not} prove that the content was not
modified. Integrity needs separate mechanisms (MACs, signatures --- later in the course).
\end{warnbox}

\section{Module 4 -- From OTP to stream ciphers}

\begin{defnbox}[title=Symmetric encryption]
Encryption of plaintext and decryption of ciphertext are done using \textbf{the same
key}. Two main types:
\begin{itemize}
  \item \textbf{Block ciphers} --- operate on fixed-size blocks (e.g.\ 128 bits).
        Example: \textbf{AES} (Advanced Encryption Standard).
  \item \textbf{Stream ciphers} --- operate one bit/byte at a time, like a
        \textbf{pseudo-OTP}. Example: \textbf{ChaCha20}.
\end{itemize}
\end{defnbox}

\begin{defnbox}[title=Symmetric cryptographic key]
A \textbf{fixed-size} input to symmetric cryptographic primitives; the size of the key
influences the level of security provided. Key properties:
\begin{itemize}
  \item \textbf{Known to both parties} --- partners must agree on the key
        \textbf{before} starting to use the primitive (requires secure generation and
        sharing protocols).
  \item \textbf{It is reused} --- the key is pre-shared once and then reused (keys do
        have a ``duration'').
  \item \textbf{It must be secret} --- revealing the key eliminates any protection
        provided by the primitive.
\end{itemize}
\end{defnbox}

\begin{rembox}[title=Difference with the OTP]
OTP key: as long as the message, used once. Symmetric key: \textbf{short and fixed}
(e.g.\ 128 or 256 bits) and \textbf{reused} for many messages. The next definitions show
how to get OTP-like encryption from such a short key.
\end{rembox}

\begin{defnbox}[title=Stream cipher (pseudo-OTP)]
A stream cipher attempts to emulate the OTP's security while solving the key-length
problem:
\begin{enumerate}
  \item \textbf{Short key} $K$ --- Alice and Bob share a small, manageable key (kept
        secret).
  \item \textbf{Key Stream Generator (KSG)} --- uses $K$ and an \textbf{Initialization
        Vector (IV)} (which does \textbf{not} need to be secret) to produce an
        arbitrarily long, \emph{pseudo-random} bit stream $S = \mathrm{stream}(K, IV)$.
  \item \textbf{Encryption} --- $C = M \oplus S$.
\end{enumerate}
To \textbf{decrypt}, the receiver feeds the KSG with the shared key and the IV (sent with
the ciphertext) to recreate the same stream $S$, then computes $C \oplus S = M$.
\end{defnbox}

\begin{propbox}[title=Security of the KSG]
We generate potentially \textbf{infinitely long randomness} (the stream $S$) from a
\textbf{finite secret} ($K$). Security depends entirely on $S$ being
\textbf{computationally indistinguishable from a truly random OTP key}: unless one knows
the key, one cannot distinguish the stream from a random string.
\end{propbox}

\begin{rembox}[title=Why does the KSG also take an IV?]
Without an IV, the same key would always produce the \textbf{same stream} $S$, and two
messages would be encrypted with the same ``pad'': XOR-ing the two ciphertexts recovers
$M_1 \oplus M_2$ (the two-time pad above), and Eve could recover the prefix of $S$ and
decode future messages. A fresh IV per message gives a \textbf{different stream for every
message} while reusing the same key.
\end{rembox}

\begin{propbox}[title=Stream ciphers: strengths and weaknesses]
\begin{center}
\begin{tabular}{p{0.32\linewidth}|p{0.58\linewidth}}
\textbf{Speed} (+) & Algorithms are linear in time and constant in space. \\
\textbf{Low error propagation} (+) & Errors in one bit do not affect subsequent
symbols. \\ \hline
\textbf{Low diffusion} ($-$) & A change in one plaintext bit only affects one ciphertext
bit. \\
\textbf{Susceptibility to modification} ($-$) & Low diffusion makes it easier to tamper
with the message (like the OTP).
\end{tabular}
\end{center}
\end{propbox}

\begin{rembox}[title=Diffusion]
\textbf{Diffusion} = how much a small change in the input spreads to the output. Stream
ciphers have none (bit $i$ of $C$ only depends on bit $i$ of $M$), which is exactly why
the bit-flipping attack works on them.
\end{rembox}

\begin{defnbox}[title=Periodicity of key stream generators]
Key stream generators which have a \textbf{finite state} are eventually
\textbf{periodic}: if the message is longer than the period, the stream
\textbf{repeats}. \\
\textbf{The attack}: if the period is identified and too small, the key stream is known
for the entire message.
\end{defnbox}

\begin{rembox}[title=Why finite state implies periodic]
A KSG with $L$ bits of internal state has at most $2^L$ different states. Each output
step moves deterministically to a next state, so after at most $2^L$ steps a state
repeats, and from then on everything repeats: a repeating stream is a \textbf{reused
pad}.
\end{rembox}

\begin{rembox}[title=Interlude on building KSGs]
The next definitions illustrate design patterns for KSGs and the challenges of building
secure ones; you are not expected to master the underlying math.
\end{rembox}

\begin{defnbox}[title=Linear Feedback Shift Register (LFSR)]
A \textbf{linear recurrence relation}: the next output bit is a linear function (XOR) of
previous bits. It extends an initial random sequence of bits into an arbitrary-length
sequence of (maybe?) random-looking bits. \\
\textbf{Advantage}: uses only \textbf{shift registers} (state storage) and \textbf{XOR
gates} (feedback), so it is extremely fast and efficient in hardware.
\end{defnbox}

\begin{exbox}[title=LFSR example]
Initial state of 4 bits $a_0, \ldots, a_3 = 1, 0, 0, 0$ and $a_n = a_{n-3} \oplus
a_{n-4}$:
\[
  a_4 = a_1 \oplus a_0 = 1,\quad a_5 = a_2 \oplus a_1 = 0,\quad
  a_6 = a_3 \oplus a_2 = 0,\quad a_7 = a_4 \oplus a_3 = 1,\ \ldots
\]
Output: $1, 0, 0, 0, 1, 0, 0, 1, 1, \ldots$
\end{exbox}

\begin{propbox}[title=Randomness of LFSRs]
\begin{itemize}
  \item \textbf{Maximal period} --- if the characteristic polynomial of the recurrence is
        \emph{primitive}, an $L$-bit register reaches the maximum period of
        $2^L - 1$ states before repeating (all non-zero states).
  \item \textbf{Distribution property} --- every non-zero state appears exactly once in
        the cycle, so the output is highly balanced: over one period it contains exactly
        $2^{L-1}$ ones and $2^{L-1} - 1$ zeros.
\end{itemize}
\end{propbox}

\begin{warnbox}[title=LFSR --- math danger]
Despite passing some statistical tests, an LFSR is purely \textbf{linear}. Techniques like
the Berlekamp-Massey algorithm (beyond the scope of the class) recover the entire LFSR
state and key stream from only a \textbf{short segment of the output}, breaking the
cipher. \textbf{LFSRs should not be used directly as key stream generators}, though they
have traditionally been used as a building block of KSGs.
\end{warnbox}

\begin{rembox}[title=Why linearity is fatal]
Linear means ``solvable with linear algebra'': each output bit gives one linear equation
on the unknown initial state. With about $2L$ output bits (easy to get in a KPA), Eve
solves the system and predicts the whole stream. ``Looks random'' (statistics) is not the
same as ``unpredictable'' (security).
\end{rembox}

\begin{exbox}[title={A5/1 and modern stream ciphers}]
\textbf{A5/1}, used historically to secure GSM mobile phone communications, combined
\textbf{multiple LFSRs} in a non-linear way to overcome the linearity flaw while staying
cheap. Despite this, mathematical weaknesses allow breaking it much faster than brute
force: it is practically broken (recall the closed-design failure of GSM, Lecture 2). \\
There are modern stream ciphers with \textbf{no known practical cryptanalysis}, e.g.\
\textbf{Trivium} (based on NLFSRs, 3 shift registers glued with XOR and AND gates) and
\textbf{Salsa20} (using different techniques).
\end{exbox}

\begin{keybox}[title=Bit security and consequences]
\begin{itemize}
  \item \textbf{Cost of breaking} --- exploiting a mathematical property or bug can break
        the scheme much faster than $2^{|K|}$ attempts, \textbf{reducing the bit
        security}.
  \item \textbf{The difficulty} --- designing mathematically secure and subtly resistant
        primitives is incredibly hard.
  \item \textbf{Conclusion} --- \textbf{never design your own} stream cipher (block
        cipher, or hash function).
  \item \textbf{Best practice} --- use \textbf{well-vetted, peer-reviewed, standard
        algorithms and their implementations} (AES, ChaCha20, SHA-3).
\end{itemize}
\end{keybox}

\section{Module 5 -- Public-key cryptography}

\begin{defnbox}[title=Key distribution problem]
Symmetric ciphers are fast, but require a \textbf{pre-shared secret key}. How can Alice
and Bob agree on a secret key if \textbf{all their communication is observed by Eve}?
\end{defnbox}

\begin{defnbox}[title=Public-key (asymmetric) cryptography]
\textbf{The paradox}: before the 1970s, it was considered impossible to communicate
securely without a pre-shared secret. \\
\textbf{The revelation}: public-key cryptography allows two parties who have
\textbf{never met} to establish a secure channel.
\end{defnbox}

\begin{defnbox}[title=Diffie-Hellman (DH) key exchange]
Alice, Bob (and Eve) know large \textbf{public parameters}: a prime number $p$ and a
generator $g$.
\begin{itemize}
  \item \textbf{Alice} chooses a \textbf{private} secret $a$, computes the
        \textbf{public value} $A = g^a \bmod p$ and sends $A$ to Bob.
  \item \textbf{Bob} chooses a \textbf{private} secret $b$, computes the
        \textbf{public value} $B = g^b \bmod p$ and sends $B$ to Alice.
\end{itemize}
Both compute the same \textbf{shared secret}:
\[
  K = B^a = (g^b)^a = g^{ab} = (g^a)^b = A^b \pmod p .
\]
Alice computes $B^a$ (she has $a$), Bob computes $A^b$ (he has $b$). Eve sees $A$, $B$,
$g$ and $p$, but cannot find $a$ or $b$ (discrete logarithm problem), so she cannot
compute $K$.
\end{defnbox}

\begin{exbox}[title={Toy example (real $p$ has 2048+ bits)}]
$p = 23$, $g = 5$. Alice picks $a = 4$: $A = 5^4 \bmod 23 = 625 \bmod 23 = 4$. Bob picks
$b = 3$: $B = 5^3 \bmod 23 = 125 \bmod 23 = 10$. \\
Alice: $K = 10^4 \bmod 23 = 18$. Bob: $K = 4^3 \bmod 23 = 64 \bmod 23 = 18$. Same key!
Eve sees $23, 5, 4, 10$ but must find $a$ with $5^a \equiv 4 \pmod{23}$.
\end{exbox}

\begin{rembox}[title=Paint analogy]
Mix a public colour with your secret colour and send the mix; each side adds its own
secret colour to what it receives: both end up with the same final colour. Un-mixing
paint (finding the secret colour) is hard, like the discrete logarithm.
\end{rembox}

\begin{defnbox}[title=Trapdoor functions]
DH is one of many ways to use public-key cryptography to establish a shared secret (a
shared key) to then do symmetric cryptography. They boil down to mathematical functions
that are \textbf{easy to compute in one direction} but \textbf{hard (high complexity) to
compute in the other direction}.
\end{defnbox}

\begin{defnbox}[title=Discrete Logarithm Problem (DLP)]
Underlying difficulty of DH: computing $g^a \bmod p$ is cheap (fast modular
exponentiation), but there is \textbf{no known fast (polynomial) algorithm} to recover
$a$ from $A = g^a \bmod p$ (and $g$, $p$).
\end{defnbox}

\begin{rembox}[title=Beyond DH]
Other ways to set a shared secret (or even to directly communicate), all involving
interesting mathematics:
\begin{itemize}
  \item \textbf{RSA} --- based on the \textbf{factoring} problem.
  \item \textbf{Elliptic Curve Cryptography (ECC)} --- the discrete logarithm problem on
        elliptic curves instead of modulo $p$ (smaller keys for the same security).
  \item \textbf{Post-quantum cryptography} --- new primitives (like \textbf{lattice
        crypto}) to resist quantum computers.
\end{itemize}
\end{rembox}

\begin{defnbox}[title=Man-In-The-Middle (MITM) attack]
Vanilla DH is vulnerable to an \textbf{active} attack: DH only guarantees key
\emph{agreement}, not \emph{authenticity} (verifying \emph{who} you are agreeing with).
\begin{enumerate}
  \item Eve intercepts $A$ from Alice and $B$ from Bob.
  \item Eve sends her own value $E_B$ to Alice: Alice computes $K_{AE}$.
  \item Eve sends her own value $E_A$ to Bob: Bob computes $K_{BE}$.
\end{enumerate}
Eve shares a key $K_{AE}$ with Alice and a separate key $K_{BE}$ with Bob: she can
\textbf{read, modify, and relay} all communication.
\end{defnbox}

\begin{rembox}[title=Why it works]
Alice believes $E_B$ comes from Bob, but nothing in the protocol proves it. Each side
does a perfectly correct DH --- just with Eve. DH is secure against a \emph{passive}
eavesdropper, not against an \emph{active} attacker (Module 2).
\end{rembox}

\begin{defnbox}[title=Identity verification problem]
DH solves \textbf{key distribution}, but creates the \textbf{identity verification
problem}.
\begin{itemize}
  \item \textbf{Solution}: DH must be paired with a mechanism (like \textbf{digital
        signatures}) to verify identity (next week).
  \item \textbf{Setup}: parties still need to trust a third party (a
        \textbf{Certificate Authority}) for public key ownership.
\end{itemize}
\end{defnbox}

\begin{keybox}[title=Summary]
\begin{itemize}
  \item Encryption $C = E_K(M)$ must be \textbf{invertible with the key} and \textbf{hard
        to invert without it}; ideally the best attack is brute force.
  \item A big key space is not enough: \textbf{frequency analysis} breaks Caesar and
        substitution ciphers. Security is measured in \textbf{bits of security}.
  \item Adversary models: \textbf{passive} $<$ \textbf{KPA} $<$ \textbf{CPA}, plus
        \textbf{side channels} on implementations.
  \item The \textbf{one-time pad} is \textbf{perfectly secret} but needs a random,
        message-length, single-use key, and is \textbf{malleable} (no integrity).
  \item \textbf{Stream ciphers} expand a short key + IV into a pseudo-random stream;
        never reuse a stream, never use plain LFSRs, \textbf{never design your own
        crypto}.
  \item \textbf{Diffie-Hellman} solves key distribution thanks to the \textbf{discrete
        logarithm problem}, but needs \textbf{authentication} against MITM.
\end{itemize}
\end{keybox}
